
\documentclass[letterpaper,11pt]{article}
%%%%%%%%%%%%%%%%%%%%%%%%%%%%%%%%%%%%%%%%%%%%%%%%%%%%%%%%%%%%%%%%%%%%%%%%%%%%%%%%%%%%%%%%%%%%%%%%%%%%%%%%%%%%%%%%%%%%%%%%%%%%%%%%%%%%%%%%%%%%%%%%%%%%%%%%%%%%%%%%%%%%%%%%%%%%%%%%%%%%%%%%%%%%%%%%%%%%%%%%%%%%%%%%%%%%%%%%%%%%%%%%%%%%%%%%%%%%%%%%%%%%%%%%%%%%
\usepackage{graphicx}
\usepackage{amsmath,amsfonts,amssymb,amsthm,float}
\usepackage{hyperref}
\usepackage[utf8]{inputenc}
\usepackage[left=2cm, right=2cm, top=2cm, bottom=2cm]{geometry}

\setcounter{MaxMatrixCols}{10}
%TCIDATA{OutputFilter=LATEX.DLL}
%TCIDATA{Version=5.50.0.2953}
%TCIDATA{<META NAME="SaveForMode" CONTENT="1">}
%TCIDATA{BibliographyScheme=BibTeX}
%TCIDATA{LastRevised=Sunday, October 04, 2026 08:44:29}
%TCIDATA{<META NAME="GraphicsSave" CONTENT="32">}
%TCIDATA{ComputeDefs=
%$\alpha =0.5449$
%$\beta =2.761\times 10^{-6}$
%$\delta =2.6486\times 10^{-6}$
%$\gamma =0.8464$
%}


\input{tcilatex}
\renewcommand{\baselinestretch}{1.15}
\setlength{\parindent}{0pt}
\setlength{\parskip}{0.5\baselineskip}
\pretolerance=2000 \tolerance=3000
\renewcommand{\abstractname}{Resumen}

\begin{document}

\title{Pr\'{a}ctica 1: Sistema de Lotka Volterra}
\author{Valadez Madera Nestor Roman $M26210087$ \\
%EndAName
Departamento de Ingenier\'{\i}a El\'{e}ctrica y Electr\'{o}nica\\
Tecnol\'{o}gico Nacional de M\'{e}xico / Instituto Tecnol\'{o}gico de Tijuana%
}
\maketitle

\noindent \textbf{Palabras clave: }Modelo Matematico; Ecuaciones
diferenciales; Simulaciones numericas; Lotka-Volterra; MATLAB.

\bigskip

\noindent Correo: \textbf{m26210087@tectijuana.edu.mx}

\noindent \noindent Carrera: \textbf{Maestr\'{\i}a en\ Ciencias de la
Ingenier\'{\i}a.}

\noindent Asignatura: \textbf{Modelado Matem\'{a}tico}

\noindent Profesor: \href{https://biomath.xyz/}{\textbf{Dr. Paul Antonio
Valle Trujillo}} (paul.vt@tijuana.tecnm.mx)

\bigskip

\section{Modelo matematico de Lotka-Volterra}

El sistema presa-depredador de Lotka-Volterra se formula mediante las
siguientes dos Ecuaciones Diferenciales Ordinarias no Lineales y de primer
orden:%
\begin{eqnarray*}
\dot{x} &=&\alpha x-\beta xy, \\
\dot{y} &=&\delta xy-\gamma y,
\end{eqnarray*}

donde los parametros $\alpha ,\beta ,\gamma ,\delta >0$ y las condiciones
iniciales $x(0),y(0)>0$. La variable de estado $x(t)$ representa a las
presas, es decir, a las liebres, mientras que la variable de estado $y(t)$
representa a los depredadores, es decir, los linces.

\bigskip

\section{Analisis de positividad}

Siguiendo el criterio general para la positividad en las soluciones de un
sistema dinamico no lineal de P. De Leenheer y D. Aeyels (2001) [Seccion 2.6
del syllabus de la asignatura], se establece el siguiente resultado:

\bigskip

\subsection{Positividad e invariancia del cuadrante no negativo.}

\textit{Sea }$\varphi (t)=(x(t),y(t))^{T}$\textit{\ una soluci\'{o}n del
sistema Lotka-Volterra con condiciones iniciales no negativas }$\varphi
(t_{0})\in 
%TCIMACRO{\U{211d} }%
%BeginExpansion
\mathbb{R}
%EndExpansion
_{+,0}^{2}$\textit{. Entonces, el cuadrante no negativo definido por }

\begin{equation*}
%TCIMACRO{\U{211d} }%
%BeginExpansion
\mathbb{R}
%EndExpansion
_{+,0}^{2}=\{\varphi \in 
%TCIMACRO{\U{211d} }%
%BeginExpansion
\mathbb{R}
%EndExpansion
_{+,0}^{2}:x,y\geq 0\}
\end{equation*}

\textit{es positivamente invariante. Por lo tanto }$\varphi \in 
%TCIMACRO{\U{211d} }%
%BeginExpansion
\mathbb{R}
%EndExpansion
_{+,0}^{2}$\textit{\ para toda }$t\geq t_{0}$\textit{.}

\bigskip

\textbf{Prueba}: Al sustituir los valores de la frontera del dominio $%
%TCIMACRO{\U{211d} }%
%BeginExpansion
\mathbb{R}
%EndExpansion
_{+,0}^{2}$, se tiene lo siguiente:

\begin{eqnarray*}
\left. \dot{x}\right\vert _{x=0} &=&\alpha (0)-\beta (0)y=0, \\
\left. \dot{y}\right\vert _{y=0} &=&\delta x(0)-\gamma (0)=0,
\end{eqnarray*}

Por lo tanto, para cada punto de la frontera evaluado en el campo vectorial
de cada ecuacion, el resultado es un valor no negativo.

\bigskip

\section{Puntos de Equilibrio}

Para calcular los puntos de equilibrio [Seccion 3.2 del syllabus de la
asignatura] es necesario igualar el campo vectorial del sistema de EDOs a
cero y resolverlo para cada variable de estado, es decir,\newline
\begin{eqnarray*}
\alpha x-\beta xy &=&0, \\
\delta xy-\gamma y &=&0,
\end{eqnarray*}

por lo tanto, al realizar el procedimiento algebraico correspondiente,

\begin{eqnarray*}
x(\alpha -\beta y) &=&0, \\
x &=&0, \\
y &=&\frac{\alpha }{\beta },
\end{eqnarray*}

al sustituir en la ecuacion de $\dot{y}$

\begin{eqnarray*}
y(\delta x-\gamma ) &=&0, \\
y &=&0, \\
x &=&\frac{\gamma }{\delta },
\end{eqnarray*}

con base en lo anterior se concluye el siguiente resultado.

\bigskip

\subsection{Resultado: Puntos de equilibrio}

\bigskip El sistema de Lotka-Volterra tiene los siguientes dos puntos de
equilibrio localizados durante el curadrante no negativo $%
%TCIMACRO{\U{211d} }%
%BeginExpansion
\mathbb{R}
%EndExpansion
_{+,0}^{2}$,

\begin{eqnarray*}
(x_{0}^{\ast },y_{0}^{\ast }) &=&(0,0), \\
(x_{1}^{\ast },y_{1}^{\ast }) &=&(\frac{\alpha }{\beta },\frac{\gamma }{%
\delta }),
\end{eqnarray*}

\bigskip

\section{Estabilidad asintotica local}

El analisis de estabilidad asintotica local de un sistema dinamico no lineal
se realiza aplicando el metodo indirecto de Lyapunov [Secciones 5.1 y 5.2.3
del syllabus de la asignatura], por lo que es necesario calcular la matriz
Jacobiana [Seccion 2.7 del syllabus de la asignatura] para linealizar el
sistema no lineal, entonces,

\begin{equation*}
\frac{\partial \varphi (t)}{\partial \varphi (x,y)}=\left[ 
\begin{array}{cc}
\frac{\partial \dot{x}}{\partial x} & \frac{\partial \dot{x}}{\partial y} \\ 
\frac{\partial \dot{y}}{\partial x} & \frac{\partial \dot{y}}{\partial y}%
\end{array}%
\right] =\left[ 
\begin{array}{cc}
\alpha -\beta y & -\beta x \\ 
\delta y & \delta x-\gamma%
\end{array}%
\right]
\end{equation*}

Para evaluar la estabilidad de cada punto de equilibrio, se evalua $J$ en
cada uno de ellos y se calculan los valores propios de la matriz resultante.
Por lo tanto, al evaluar $(x_{0}^{\ast },y_{0}^{\ast })$ que representa el
orginen, se obtiene el siguiente resultado:

\begin{equation*}
\left. J\right\vert _{(x_{0}^{\ast },y_{0}^{\ast })}=\left[ 
\begin{array}{cc}
\alpha & 0 \\ 
0 & -\gamma%
\end{array}%
\right] ,
\end{equation*}

Los valores propios de una matriz diagonal son cada uno de los elementos de
su diagonal, es decir,

\begin{eqnarray*}
\lambda _{1} &=&\alpha , \\
\lambda _{2} &=&-\gamma ,
\end{eqnarray*}

\bigskip

Ahora, al evaluar el punto de equilibrio $(x_{1}^{\ast },y_{1}^{\ast })=(%
\frac{\alpha }{\beta },\frac{\gamma }{\delta })$ en $J$, se tiene lo
siguiente:

\begin{equation*}
\left. J\right\vert _{(x_{1}^{\ast },y_{1}^{\ast })}=\left[ 
\begin{array}{cc}
0 & -\frac{\beta \gamma }{\delta } \\ 
\frac{\alpha \delta }{\beta } & 0%
\end{array}%
\right] ,
\end{equation*}

\subsubsection{\textbf{Tarea:} Que se concluye con respecto \ a la
estabilidad del punto de equilibrio $(x_{0}^{\ast },y_{0}^{\ast })$?}

\bigskip Utilizando los valores para $\alpha ,\beta ,\gamma ,\delta $
obtenidos en MATLAB, se resuelve el punto de equilibrio $(x_{0}^{\ast
},y_{0}^{\ast })$ de la siguiente manera:

$\alpha =0.5449$

$\beta =2.761\times 10^{-6}$

$\delta =2.6486\times 10^{-6}$

$\gamma =0.8464$

Solution is: $0.544\,9,-0.846\,4\allowbreak $%
\begin{eqnarray*}
\det (%
\begin{array}{cc}
\alpha & 0 \\ 
0 & -\gamma%
\end{array}%
-%
\begin{array}{cc}
\lambda & 0 \\ 
0 & \lambda%
\end{array}%
) &=&0 \\
\allowbreak \lambda _{1} &=&\alpha \\
\lambda _{2} &=&-\gamma
\end{eqnarray*}

El metodo indirecto de Lyapunnov expresa que: \textit{El origen de }$\dot{x}$%
\textit{\ es localmente asintoticamente inestable si }$\func{Re}\lambda
_{i}>0$\textit{\ \ para uno o m\'{a}s valores propios de }$J$\textit{.} Uno
de los valores propios es $\lambda _{1}=-\alpha $, concluyendo que el punto
de equilibrio $(x_{0}^{\ast },y_{0}^{\ast })$ es inestable.

\subsubsection{\textbf{Tarea}: Cuales son los valores propios de la matriz
Jacobiana resultante al evaluar el punto de equilibrio $(x_{1}^{\ast
},y_{1}^{\ast })$?}

\begin{eqnarray*}
\det (%
\begin{array}{cc}
0 & -\frac{\beta \gamma }{\delta } \\ 
\frac{\alpha \delta }{\beta } & 0%
\end{array}%
-%
\begin{array}{cc}
\lambda & 0 \\ 
0 & \lambda%
\end{array}%
) &=&\allowbreak \lambda ^{2}+0.461\,2 \\
\lambda _{1} &=&-0.679\,12i \\
\lambda _{2} &=&0.679\,12i
\end{eqnarray*}

\subsubsection{Tarea: Que se concluye con respecto \ a la estabilidad del
punto de equilibrio $(x_{1}^{\ast },y_{1}^{\ast })$?}

\textit{El teorema del m\'{e}todo indirecto de Lyapunov no considera los
casos en que }$\func{Re}\lambda _{i}\leq 0$\textit{\ o }$\func{Re}\lambda
_{i}=0$\textit{, en estos el teorema no se puede aplicar para concluir sobre
la estabilidad local del punto de equilibrio de }$\dot{x}=f(x)$\textit{.}
Sin embargo, es posible realizar conclusiones sobre la estabilidad y no
sobre la geometr\'{\i}a detallada de las trayectorias. Acorde al syllabus,
los casos en los que $\func{Re}\lambda _{i}=0$ son considerados\textit{\
Casos Marginales }y en este caso se les conoce como \textit{centros}. Es
decir, la trayectoria al rededor del equilibrio presentara oscilaciones
sostenidas.

\bigskip 

\section{Normalizaci\'{o}n (escalamiento) del sistema}

El sistema din\'{a}mico de Lotka-Volterra

\begin{eqnarray*}
\dot{x} &=&\alpha x-\beta xy, \\
\dot{y} &=&\delta xy-\gamma y,
\end{eqnarray*}%
\newline
se normaliza de la siguiente forma, se realiza el cambio de variable:\newline

\begin{equation*}
x=\frac{x_{n}}{x_{\max }}\text{ y }y=\frac{y_{n}}{y_{\max }},
\end{equation*}

por lo tanto

\begin{equation*}
\frac{dx}{dt}=\frac{1}{x_{\max }}\frac{dx_{n}}{dt}\text{ y }\frac{dy}{dt}=%
\frac{1}{y_{\max }}\frac{dy_{n}}{dt},
\end{equation*}

al sustituir

\begin{eqnarray*}
\frac{1}{x_{\max }}\frac{dx_{n}}{dt} &=&\alpha \frac{x_{n}}{x_{\max }}-\beta 
\frac{x_{n}}{x_{\max }}\frac{y_{n}}{y_{\max }}=\alpha x_{n}-\frac{\beta }{%
y_{\max }}x_{n}y_{n}, \\
\frac{1}{y_{\max }}\frac{dy_{n}}{dt} &=&\delta \frac{x_{n}}{x_{\max }}\frac{%
y_{n}}{y_{\max }}-\gamma \frac{y_{n}}{y_{\max }}=\frac{\delta }{x_{\max }}%
x_{n}y_{n}-\gamma y_{n},
\end{eqnarray*}

finalmente:

\begin{equation*}
\beta _{n}=\frac{\beta }{y_{\max }}\text{ y }\delta _{n}=\frac{\delta }{%
x_{\max }},
\end{equation*}

El parametro original se obtiene al realizar

\begin{eqnarray*}
\beta  &=&\beta _{n}y_{\max }, \\
\delta  &=&\delta _{n}x_{\max },
\end{eqnarray*}

\section{Discusion de resultados y simulaciones}

Durante la pr\'{a}ctica se busca aplicar el sistema de Lotka-Volterra para
analizar sus caracteristicas como modelo matematico.

En esta se realizaron simulaciones numericas en MATLAB, Simulink y se
realizo el an\'{a}lisis matem\'{a}atico en Scientific Workplace.

\bigskip

\subsubsection{Scientific Workplace}

Para el analisis matem\'{a}tico se llegaron a las siguientes conclusiones:

1.- El sistema permanecera en $%
%TCIMACRO{\U{211d} }%
%BeginExpansion
\mathbb{R}
%EndExpansion
_{+,0}^{2}$ para toda $\{x,y\}\geq 0$.

2.- El sistema cuenta con dos puntos de equilibrio $\{x_{e0},y_{e0}\}=\{0,0\}
$ y $\{x_{e1},y_{e1}\}=\{\alpha ,-\gamma \}$ calculados por medio de la
linealizaci\'{o}n de la forma $\dot{x}=Ax$.

\bigskip

De los puntos de equlibrio se puede concluir lo siguiente:

2.1.- El metodo indirecto de Lyapunnov expresa que: \textit{El origen de }$%
\dot{x}$\textit{\ es localmente asintoticamente inestable si }$\func{Re}%
\lambda _{i}>0$\textit{\ \ para uno o m\'{a}s valores propios de }$J$\textit{%
.} Uno de los valores propios es $\lambda _{1}=\alpha $, concluyendo que el
punto de equilibrio $(x_{0}^{\ast },y_{0}^{\ast })$ es inestable. Esta din%
\'{a}mica se comporta como un repulsor.

2.2.- \textit{El teorema del m\'{e}todo indirecto de Lyapunov no considera
los casos en que }$\func{Re}\lambda _{i}\leq 0$\textit{\ o }$\func{Re}%
\lambda _{i}=0$\textit{, en estos el teorema no se puede aplicar para
concluir sobre la estabilidad local del punto de equilibrio de }$\dot{x}%
=f(x) $\textit{.} No es posible realizar conclusiones sobre la estabilidad
pero s\'{\i} sobre la geometr\'{\i}a detallada de las trayectorias. Acorde
al syllabus, los casos en los que $\func{Re}\lambda _{i}=0$ son considerados%
\textit{\ Casos Marginales }y en este caso se les conoce como \textit{centros%
}. Es decir, la trayectoria al rededor del equilibrio presentara
oscilaciones sostenidas.

\bigskip

\subsubsection{MATLAB y Simulink}

En MATLAB se realizo la aproximaci\'{o}n de los parametros $\alpha $,$\beta $%
,$\gamma $ y $\delta $ para los datos crudos y la simulaci\'{o}m de la ecuaci%
\'{o}n cuando se normalizan y/o suavizan los datos. Adiconalmente se
realizaron simulaciones numericas sobre los parametros obtenidos utilizando
los datos crudos las cuales utilizar diferentes odes desde la herramienta de
Simulink.

La primera funci\'{o}n realizada para la obtenci\'{o}n de parametros fue la
funci\'{o}n mdl, la cual requiere unos parametros preliminares para la
ooptimizaci\'{o}n de estos, la otra funci\'{o}n es la funcion LotkaVolterra
que utiliza los parametros obtenidos para el modelado directo de la ecuaci%
\'{o}n, ambas funciones utilizan el m\'{e}todo de Euler-Heun para realizar
la aproximaci\'{o}n matem\'{a}tica de la funci\'{o}n. Las funciones plotsys,
plotfittign y plotphase se utilizar\'{o}n para gr\'{a}ficar las funciones
obtenidas en diferentes puntos de la pr\'{a}ctica.

Para el computo de la soluci\'{o}n se utilizaron odes para funciones stiff
(ode45, ode23, ode113), odes con paso de integraci\'{o}n fijo (ode1, ode2,
ode3, ode4 y ode5) y funciones anonimas (ode78, ode89). Las funciones con
paso de integraci\'{o}n fijo ten\'{\i}an una menor cantidad de oscilaciones
en general que las funciones anonimas y las odes para funciones stiff.

\bigskip

En la graficaci\'{o}n de la estabilidad de los puntos de equilibrio podemos
encontrar varios casos:

1.- Caso en la vecindad del punto $(x_{0}^{\ast },y_{0}^{\ast })$ y alejado
de $(x_{1}^{\ast },y_{1}^{\ast })$: En este caso una forma de elipsoide
bastanre deformada, si bien el punto $(x_{0}^{\ast },y_{0}^{\ast })$ \ es un
punto de inestabilidad, los valores de la condicion inicial utilizada se
encuentran entre $(x_{0}^{\ast },y_{0}^{\ast })$ y $(x_{1}^{\ast
},y_{1}^{\ast })$ causando que la funcion fase tenga la forma de un
elipsoide bastante deformado al grado de parecer un tri\'{a}ngulo rectangulo.

2.- Caso en la vecindad del punto $(x_{1}^{\ast },y_{1}^{\ast })$ y no tan
alejado de $(x_{0}^{\ast },y_{0}^{\ast })$. En este caso, las condiciones
iniciales se encutran cerca del punto de equilibrio $(x_{1}^{\ast
},y_{1}^{\ast })$ por un factor de ($\pm 50\%$ $(x_{1}^{\ast },y_{1}^{\ast })
$), y mientras mas se disminuya este factor la forma del centroide se
volvera cada vez mas pronunciada, en este segundo caso por ejemplo, la forma
de la funci\'{o}n fase pasa a tener una forma similar a la de un huevo.

3.- Caso en la vecindad del punto $(x_{1}^{\ast },y_{1}^{\ast })$ y alejado
de $(x_{0}^{\ast },y_{0}^{\ast })$. En este caso, las condiciones iniciales
se encuentran cerca del punto de equilibrio ($\pm 5\%$ $(x_{1}^{\ast
},y_{1}^{\ast })$) la funci\'{o}n pasa a ser un ovalo con de un radio
reducido alrededor del punto de equilibrio.

4.- Caso alejado de $(x_{0}^{\ast },y_{0}^{\ast })$ y de $(x_{1}^{\ast
},y_{1}^{\ast })$: En este caso, la funcion fase nuevamente vuelve a
deformase mientras mas se aleja del equilibrio $(x_{1}^{\ast },y_{1}^{\ast
}) $, por ejemplo para un factor de incremento de $200\%$ $(x_{1}^{\ast
},y_{1}^{\ast })$ la forma de la funci\'{o}n fase vuelve a parecer la de un
triangulo rectangulo y continua deformandose al alejarse m\'{a}.

\end{document}
